\scanpage{146}
Sur l'ens des \deleted{\unsure{structu|structures}} classes d'isomorphie de structures de Torelli directes, le groupe
\[
\deleted{\unsure{SL(2,}}
\]
\[
SL(2,\mathbb Z)=\left\{\begin{pmatrix}a&b\\c&d\end{pmatrix}\;\middle|\;
\begin{aligned}a,b,c,d&\in\mathbb Z\\ad-bc&=1\end{aligned}\right\} \tag*{(9)}
\]
opère à droite, par
\[
t\unsure{\mathrel{\vdash\!\dashv}|\longmapsto}t\circ g,
\]
on va expliciter cette opération, en tant qu'opération sur $\mathcal D^+$, et $\unsure{\overline{\mathrm m}|=}$ sur $X_{\mathcal D^+}$ (vu que l'on a une opération de $SL(\unsure{2|1},\mathbb Z)$ sur la situation modulaire).

On voit qu'à une structure de Torelli $(E,(z_1,z_2))$, $g$ associe la structure $(E,(z'_1,z'_2))$, avec
\[
\begin{aligned}
z'_1&=az_1+cz_2\\
z'_2&=bz_1+dz_2
\end{aligned}
\]
d'où, si $\tau,\tau'$ sont les invariants correspondants
\[
-\qquad \tau=z_2/z_1,\quad\tau'=z'_2/z'_1
\]
la formule
\[
\tau'=\frac{d\tau+b}{\deleted{\unsure{b}}c\tau+a}\qquad
\left(\tau'=\tau\circ g,\ g=\begin{pmatrix}a&b\\c&d\end{pmatrix}\right) \tag*{(10)}
\]
Si on remplace l'opération envisagée par l'opération à gauche correspondante, \deleted{\unsure{on|---}} $t\longmapsto t\circ g^{-1}$, notée aussi $\tau\longmapsto\tau\circ g^{-1}$, on trouve,
% archive sheet number: 145
